\clearpage\SourceLeaf{186}%
\pdfbookmark[1]{Neue Aussicht auf eine Erweiterung des Gebiets der Astronomie}{aussicht1813}
\fancyhead[L]{\small NEUE AUSSICHT FÜR DIE ASTRONOMIE}
\begin{center}
{\small NEUE AUSSICHT AUF EINE ERWEITERUNG\par
DES GEBIETS DER ASTRONOMIE.\par}
\medskip\hrule\medskip
{\small Neues Göttingisches Taschenbuch zum Nutzen und Vergnügen für das Jahr 1813.\par
Göttingen, bei Heinrich Dieterich.\par}
\medskip\hrule
\end{center}\medskip
So weit die Geschichte zurückreicht, war die Erde stets ein unruhiger Schauplatz feindseliger Leidenschaften. Ein Volk verdrängt das andere, um früher oder später selbst wieder einem anderen Platz zu machen. Ein Reich erhebt sich auf den Trümmern eines anderen. Vergänglichkeit ist das Los aller menschlichen Schöpfungen, an deren Segen oder Fluch die Generationen teilhaben, die der Strom der Zeit gerade an ihnen vorüberführt. Dieser ewige Wechsel aller Dinge auf der Erde ist demütigend und bedrückend. Erhebend und aufheiternd ist dagegen der Anblick, wie der menschliche Geist auf dem Gebiet der Wahrheit unablässig voranschreitet, und nirgends mehr als in der erhabensten der Naturwissenschaften, der Astronomie.

Seit zwei Jahrtausenden sehen wir, wie sich die schöne Frucht der Beobachtung und des tiefen Nachdenkens mitten in den Stürmen der Zeiten entwickelt und wächst, gepflegt von den liebevollen Händen derer, die sich der Astronomie verschrieben haben. Ein Jahrhundert gibt das kostbare Vermächtnis an das nächste weiter, damit es dieses weiter vermehrt. Immer mehr von ihren wunderbaren Geheimnissen enthüllt die Himmlische. Die Sterne sind gezählt. Gemessen ist die scheinbar unermessliche Entfernung der Sonne und der Himmelskörper, die sie umkreisen. Ihre Bahnen sind genau bestimmt. Eine einzige, erhaben einfache Grundkraft erklärt ihre verwickelten Bewegungen und alle vielfältigen Erscheinungen, die sich daraus ergeben, vollständig. Gemessen,
\expandafter\def\csname HistoricalNoteText13\endcsname{Die Beobachtungen von \textsc{Piazzi}, \textsc{Calandrelli} und \textsc{Brinkley} ergeben zwar eine kleine jährliche Parallaxe für den hellen Stern in der Leier. Doch die Beobachtungen sind noch mit Zweifeln behaftet, und die Ergebnisse stimmen nicht überein. Aus \textsc{Bradleys} Beobachtungen ergibt sich für denselben Stern zudem überhaupt keine merkliche Parallaxe.}%
\expandafter\def\csname HistoricalNoteText14\endcsname{Mithilfe der Durchgänge von 1761 und 1769.}%
\expandafter\def\csname HistoricalNoteText15\endcsname{Mithilfe der einen kleinen Ungleichheit in der Bewegung des Mondes, die von der Sonnenparallaxe abhängt.}%
\SourceLeaf{187}%
ja gewogen, sind die Erde, die Sonne, die Planeten und die Monde. Gemessen ist die Geschwindigkeit des Lichtstrahls. Verschwunden sind die Wahnvorstellungen, die einst die Völker ängstigten. Auch dem Irdischen zu dienen verschmäht die Himmlische nicht. Sie teilt und zählt dem Menschen seine Zeiträume, vermisst und kartiert ihm Länder und Kontinente und leitet seine Schiffe sicher über den Ozean.

Unser Sonnensystem ist der eigentliche Schauplatz unserer bisherigen Astronomie. Noch viele Jahrhunderte lang wird es unerschöpflichen Stoff für neue Forschungen bieten. Doch im Weltall nimmt unser Sonnensystem einen noch viel kleineren Platz ein als die Erde in unserem Sonnensystem. Es ist ein Tropfen im Ozean. Das Heer der Fixsterne ist für uns bisher kaum mehr als eine Anzahl fester Punkte, auf die wir unsere Beobachtungen beziehen. Von ihren unermesslichen Entfernungen haben wir nur eine Ahnung. Verglichen damit verschwindet fast die längste Grundlinie, mit der wir sie zu messen versuchen, der Durchmesser der jährlichen Erdbahn\HistoricalNote{13}. Geben wir deshalb die Hoffnung nicht auf, etwas über diese Entfernungen zu erfahren. Auch die Entfernung der Sonne hat man mit dem kleinen Erddurchmesser nicht unmittelbar gemessen. Unter günstigen Umständen wählte man dafür die nähere Venus{\renewcommand{\thefootnote}{**)}\HistoricalNote{14}}, und später wusste man mit Scharfsinn die schon länger bekannte Entfernung des Mondes als Maßstab für die Entfernung der Sonne zu nutzen{\renewcommand{\thefootnote}{***)}\HistoricalNote{15}}. Vielleicht sind uns gerade die hellsten Fixsterne nicht am nächsten. Was bei ihnen bisher noch nicht ganz gelungen ist, gelingt vielleicht eher bei einem unscheinbaren Stern von viel geringerer Helligkeit.

Tatsächlich gibt es gewichtige Gründe für diese Vermutung. Die sogenannten Fixsterne, zu denen auch unsere Sonne gehört, sind im strengen Sinn des Wortes nicht fest. Selbst wenn sie einmal alle zugleich in Ruhe gewesen wären, hätten sie wegen ihrer gegenseitigen Einwirkungen nicht in Ruhe bleiben können. Auch in diesem zahllosen Heer von Fixsternen ist überall Bewegung. Hier muss man räumliche Entfernungen ebenso wie Zeiträume in großen Maßstäben messen.
\SourceLeaf{188}%
Ein Jahrhundert verändert ihre gegenseitigen Stellungen weniger, als ein Tag die gegenseitigen Stellungen der Planeten verändert. Doch mit der Genauigkeit unserer heutigen Beobachtungen lassen sich diese Stellungsänderungen über fünfzig oder hundert Jahre bei einer beträchtlichen Zahl von Sternen bereits erkennen. Unter gleichen Umständen fallen diese Änderungen bei den \emph{nächsten} Sternen am meisten auf. Merkwürdigerweise sind sie bei einigen schwächeren Sternen beträchtlicher als bei den Sternen erster Größe. \emph{Solche} Sterne, von denen man vielleicht noch mehrere entdecken wird, sollte man besonders sorgfältig auf ihre jährliche Parallaxe hin beobachten, also darauf, wie sich ihre Stellung je nach dem Standort der Erde auf ihrer Bahn ändert.

Ein Stern verdient in dieser Hinsicht ganz besondere Aufmerksamkeit. Es ist der Stern sechster Größe im Sternbild Schwan, den \textsc{Flamstead} mit der Zahl 61 bezeichnet hat. Von allen Sternen, bei denen man bisher eine Verschiebung wahrgenommen hat, zeigt dieser die größte, nämlich 8 Bogenminuten und 52 Bogensekunden in hundert Jahren. Eine so beträchtliche Änderung seiner Stellung zu uns setzt entweder eine sehr schnelle tatsächliche Bewegung im Weltraum oder eine verhältnismäßig geringere Entfernung voraus. Diese wichtige Beobachtung verdanken wir Herrn \textsc{Bessel} in Königsberg. Er kam darauf, als er \textsc{Bradleys} Beobachtungen dieses Sterns mit Beobachtungen aus den letzten Jahren des vorigen Jahrhunderts verglich.

Ein anderer Umstand macht die beobachtete schnelle Bewegung dieses Sterns noch viel interessanter. Der Stern ist eigentlich ein sogenannter \emph{Doppelstern}. Durch ein gutes Fernrohr sieht man, dass er aus zwei Sternen besteht, die sehr nahe beieinander erscheinen (ihr Abstand beträgt 20 Bogensekunden, und ihre Helligkeit unterscheidet sich kaum). Solche Doppelsterne gibt es am Himmel in großer Zahl. Die Astronomen waren bisher stets der Meinung, dass zwei Sterne uns nur zufällig als Doppelstern erscheinen, weil sie mit uns beinahe auf einer einzigen Geraden liegen. Ihre Entfernungen von uns können dabei sehr verschieden sein. Vor wenigen Jahren hat Herr \textsc{Herschel} als Erster diese Ansicht bestritten und behauptet, die Sterne solcher Paare stünden tatsächlich nahe beieinander im Raum und bildeten ein eigenes System. Die Gründe, die \textsc{Herschel} für diese Ansicht anführt, sind freilich nicht entscheidend. Doch schon an sich ist es äußerst unwahrscheinlich,
\SourceLeaf{189}%
dass der Zufall so viele Doppelsterne mit so geringem Abstand gebildet haben sollte. Viel eher kann man darauf wetten, dass zwischen ihnen eine wirkliche Verbindung besteht. Unser Stern im Schwan fällt nun stark zugunsten von \textsc{Herschels} Ansicht ins Gewicht. Besteht zwischen den beiden Sternen keine weitere Verbindung, und ändern sie ihren Platz im Weltraum oder wir den unseren, so ist mit größter Wahrscheinlichkeit zu erwarten, dass sie uns bald nicht mehr als Doppelstern erscheinen. Am Himmel erscheinen zum Beispiel oft kleine Sterne nahe bei Jupiter. Wenn der Planet weiterzieht, scheint er die Sterne zurückzulassen. Nur die Trabanten, die ihm wirklich nahe sind, bleiben ihm stets treu. Unser Doppelstern ist aber in 50 Jahren am Firmament dreizehnmal so weit vorgerückt, wie der Abstand zwischen den beiden Sternen beträgt, und der Abstand selbst ist fast unverändert geblieben. Nur stehen die Sterne jetzt weniger schräg zueinander (1755 stand der eine nordöstlich vom anderen, jetzt fast genau östlich). Es bleibt also kaum ein Zweifel, dass wir in ihnen eine neue Besonderheit im Aufbau der Welt kennengelernt haben, ein Sonnensystem, in dem zwei Sonnen anstelle einer einzigen herrschen.

Eine reine, einseitige Herrschaft gibt es in der Welt der Körper nicht. Was herrscht, erfährt immer mehr oder weniger die Rückwirkung dessen, worüber es herrscht. Jene beiden Sonnen werden wechselseitig übereinander herrschen. Die zweite wird nicht um die erste laufen, die erste nicht um die zweite. Beide werden ihren gemeinsamen Schwerpunkt umlaufen. Je näher sie einander sind und je stärker sie aufeinander einwirken, desto schneller werden sie einen solchen wechselseitigen Umlauf vollenden. Wie weit sie voneinander entfernt sind, werden wir beurteilen können, wenn sich die Hoffnung als begründet erweist, dass wir ihre Entfernung von uns noch messen können. Und wie lange dieser Umlauf dauert, der vielleicht nur wenige Jahrhunderte währt, wird die Zeit lehren. Dann wird man in der Kenntnis des Sternenhimmels einen großen Schritt weiter sein. Man wird jene Sterne \emph{wiegen} können, zumindest ihre Gesamtmasse, auch wenn das bei den einzelnen Sternen nicht gelingt. Dem ersten Schritt werden weitere folgen. Aber wer zählt die Jahrtausende, die erst vergehen müssen, bis unsere Kenntnisse vom Sternenhimmel auch nur ein Schatten dessen werden können, was wir längst über \emph{unser} so großes und so kleines Sonnensystem wissen? Und wie viel Neues bleibt auch darin immer noch zu lernen.
\begin{center}\rule{0.24\linewidth}{0.4pt}\end{center}
